\clearpage
\chapter*{LAMPIRAN F. TABEL RANCANGAN LENGKAP}
\addcontentsline{loa}{appendix}{Lampiran F. Tabel Rancangan Lengkap}

Lampiran ini memuat rincian tabel rancangan yang dipisahkan dari pembahasan utama.
Setiap tabel berikut disusun sebagai satu kesatuan agar dapat dibaca utuh tanpa
potongan halaman yang tidak diperlukan.

Rincian skema dan konfigurasi dirujuk pada Tabel~\ref{tab:rancangan-schema-pipeline-profile-nodes},
Tabel~\ref{tab:rancangan-schema-rag-node-runs}, Tabel~\ref{tab:rancangan-node-event-envelope},
Tabel~\ref{tab:rancangan-node-message-chunker}, Tabel~\ref{tab:rancangan-canonical-chunker},
Tabel~\ref{tab:rancangan-canonical-kg}, dan Tabel~\ref{tab:rancangan-config-indexer-dense}.

\begingroup
\scriptsize
\setlength{\tabcolsep}{2pt}
\setcounter{table}{0}
\renewcommand{\thetable}{F.\arabic{table}}

\begin{longtable}{@{}p{3.8cm}p{2.8cm}p{2.5cm}p{\dimexpr\textwidth-3.8cm-2.8cm-2.5cm-6\tabcolsep\relax}@{}}
\caption{Struktur tabel \texttt{pipeline\_profile\_nodes}: urutan \emph{Node} dan \emph{Plugin}}\label{tab:rancangan-schema-pipeline-profile-nodes}\\
\toprule \emph{Field} & Tabel/Tipe & Kunci/NULL & Fungsi\\ \midrule
\endfirsthead
\multicolumn{4}{l}{\footnotesize\emph{(lanjutan)}}\\ \toprule \emph{Field} & Tabel/Tipe & Kunci/NULL & Fungsi\\ \midrule \endhead
\bottomrule \endfoot
\code{id} & \code{uuid} & PK, NOT NULL & Identitas baris \emph{node}.\\
\code{profile_id} & \code{uuid} & FK, NOT NULL & \emph{Profile} pemilik, \emph{node} dihapus bersama \emph{profile}.\\
\code{node_id} & \code{text} & UNIQUE/\emph{profile} & Identitas \emph{node} di dalam \emph{profile}.\\
\code{component_type} & enum \code{component_type} & NOT NULL & \code{parser}, \code{chunker}, \code{indexer}, \code{kg}, \code{retrieval_orchestrator}, \code{xai_intent}, \code{xai_retrieval}, \code{xai_reranking}, \code{scraper}.\\
\code{plugin_id} & \code{uuid} & FK, NOT NULL & \emph{Plugin} yang dipilih.\\
\code{sort_order} & \code{integer} & UNIQUE/\emph{profile} & Posisi \emph{node}, satu posisi tidak boleh ganda.\\
\code{config_json} & \code{json} & NOT NULL & Konfigurasi efektif yang lolos JSON \emph{Schema}.\\
\end{longtable}

\begin{longtable}{@{}p{3.8cm}p{2.8cm}p{2.5cm}p{\dimexpr\textwidth-3.8cm-2.8cm-2.5cm-6\tabcolsep\relax}@{}}
\caption{Struktur tabel \texttt{rag\_node\_runs}: status dan percobaan \emph{Node}}\label{tab:rancangan-schema-rag-node-runs}\\
\toprule \emph{Field} & Tabel/Tipe & Kunci/NULL & Fungsi\\ \midrule
\endfirsthead
\multicolumn{4}{l}{\footnotesize\emph{(lanjutan)}}\\ \toprule \emph{Field} & Tabel/Tipe & Kunci/NULL & Fungsi\\ \midrule \endhead
\bottomrule \endfoot
\code{id} & \code{uuid} & PK, NOT NULL & Identitas satu eksekusi \emph{node}.\\
\code{run_id} & \code{uuid} & FK CASCADE & \emph{Pipeline run} pemilik.\\
\code{component_type} & enum \code{rag_component_type} & NOT NULL & \code{chunker}, \code{indexer}, \code{kg}, \code{parser}, \code{scraper}.\\
\code{plugin_name} & \code{text} & NULL & Nama \emph{plugin} yang dijalankan.\\
\code{plugin_version} & \code{text} & NULL & Versi \emph{plugin} yang dijalankan.\\
\code{status} & enum \code{rag_run_status} & NOT NULL & \code{queued}, \code{running}, \code{succeeded}, \code{failed}, \code{cancelled}.\\
\code{attempt} & \code{integer} & NOT NULL & Nomor percobaan \emph{node}.\\
\code{pipeline_attempt} & \code{integer} & NOT NULL & Nomor percobaan \emph{pipeline}.\\
\code{artifact_object_key} & \code{text} & NULL & \emph{Artifact} keluaran \emph{node}.\\
\code{cache_hit} & \code{boolean} & NULL & Penanda keluaran dipulihkan dari \emph{artifact cache}.\\
\code{cache_key} & \code{text} & NULL & Kunci deterministik untuk pencarian \emph{artifact cache}.\\
\code{error_json} & \code{jsonb} & NULL & Kegagalan lokal \emph{node}.\\
\code{created_at} & \code{timestamptz} & NOT NULL, \code{now()} & Waktu \emph{node} \emph{run} dibuat.\\
\code{started_at} & \code{timestamptz} & NULL & Waktu \emph{node} mulai diproses.\\
\code{finished_at} & \code{timestamptz} & NULL & Waktu \emph{node} selesai.\\
\code{updated_at} & \code{timestamptz} & NOT NULL, \code{now()} & Waktu status \emph{node} terakhir diubah.\\
\end{longtable}

\begin{longtable}{@{}p{3.5cm}p{3.0cm}p{\dimexpr\textwidth-3.5cm-3.0cm-4\tabcolsep\relax}@{}}
\caption{Skema \texttt{NodeEventEnvelope}}\label{tab:rancangan-node-event-envelope}\\
\toprule Struktur/\emph{Field} & Tipe atau Nilai & Fungsi\\ \midrule
\endfirsthead
\multicolumn{3}{l}{\footnotesize\emph{(lanjutan)}}\\ \toprule Struktur/\emph{Field} & Tipe atau Nilai & Fungsi\\ \midrule \endhead
\bottomrule \endfoot
\code{event_id} & \emph{string} non-empty & Identitas unik \emph{event}.\\
\code{schema_version} & literal \code{"2"} & Versi kontrak \emph{message}.\\
\code{run_id} & \emph{string} non-empty & \emph{Pipeline run} pemilik.\\
\code{requested_at} & datetime & Waktu permintaan.\\
\code{request_id} & \emph{string} non-empty & Korelasi permintaan awal.\\
\code{node_id} & \emph{string} non-empty & \emph{Node profile} yang dituju.\\
\code{attempt} & \emph{integer} $\geq 1$ & Nomor percobaan \emph{node}.\\
\end{longtable}

\begin{longtable}{@{}p{3.5cm}p{3.0cm}p{\dimexpr\textwidth-3.5cm-3.0cm-4\tabcolsep\relax}@{}}
\caption{Skema pesan permintaan \emph{chunker}}\label{tab:rancangan-node-message-chunker}\\
\toprule \emph{Field} & Tipe & Fungsi\\ \midrule
\endfirsthead
\multicolumn{3}{l}{\footnotesize\emph{(lanjutan)}}\\ \toprule \emph{Field} & Tipe & Fungsi\\ \midrule \endhead
\bottomrule \endfoot
\code{idempotency_key} & \emph{string} & Kunci deduplikasi untuk satu permintaan \emph{chunker}.\\
\code{event} & \code{NodeEventEnvelope} & Identitas \emph{run}, \emph{node} \emph{chunker}, \emph{attempt}, dan korelasi.\\
\code{profile_id} & \emph{string}, NULL & \emph{Profile} \emph{target} ketika pesan memerlukan identitas eksplisit.\\
\code{document_artifact} & \code{ArtifactReference}, NULL & Referensi \code{canonical-document}.\\
\code{document_object_key} & \emph{string}, NULL & Referensi kompatibilitas, minimal satu sumber dokumen harus tersedia.\\
\end{longtable}

\begin{longtable}{@{}>{\RaggedRight\arraybackslash}p{4.0cm}>{\RaggedRight\arraybackslash}p{3.0cm}>{\RaggedRight\arraybackslash}p{2.4cm}>{\RaggedRight\arraybackslash}p{\dimexpr\textwidth-4.0cm-3.0cm-2.4cm-6\tabcolsep\relax}@{}}
\caption{Kontrak \emph{canonical artifact} \emph{chunker}}\label{tab:rancangan-canonical-chunker}\\
\toprule Kontrak/Elemen & \emph{Field} & Tipe & Fungsi\\ \midrule
\endfirsthead
\multicolumn{4}{l}{\footnotesize\emph{(lanjutan)}}\\ \toprule Kontrak/Elemen & \emph{Field} & Tipe & Fungsi\\ \midrule \endhead
\bottomrule \endfoot
\code{CanonicalChunks} & \code{document_id} & \emph{string} & Identitas dokumen asal kumpulan \emph{chunk}.\\
\code{CanonicalChunks} & \code{chunker} & \emph{PluginRef} & Identitas \emph{plugin} \emph{chunker} pembentuk.\\
\code{CanonicalChunks} & \code{chunks} & \emph{array of} \code{Canonical}\allowbreak\code{Chunk} & Daftar elemen \code{CanonicalChunk} hasil pemotongan.\\
\code{CanonicalChunks} & \code{metadata} & \emph{object} & \emph{Metadata} kumpulan \emph{chunk}.\\
\code{CanonicalChunk} & \code{chunk_id} & \emph{string} & Identitas \emph{chunk}.\\
\code{CanonicalChunk} & \code{kind} & \emph{enum} & Kanal penggunaan \emph{chunk}, yaitu \emph{retrieval}, \emph{context}, atau \emph{both}.\\
\code{CanonicalChunk} & \code{text} & \emph{string} & Teks \emph{chunk}.\\
\code{CanonicalChunk} & \code{metadata} & \emph{object} & \emph{Metadata} \emph{chunk}.\\
\code{CanonicalChunk} & \code{source_blocks} & \emph{array of} \code{string} & Referensi ID \emph{block} sumber \emph{chunk}.\\
\code{relationships} & \code{parent_id} & \emph{string}, NULL & Referensi \emph{parent} \emph{chunk} jika tersedia.\\
\code{relationships} & \code{previous_id} & \emph{string}, NULL & Referensi \emph{chunk} sebelumnya jika tersedia.\\
\code{relationships} & \code{next_id} & \emph{string}, NULL & Referensi \emph{chunk} berikutnya jika tersedia.\\
\end{longtable}

\clearpage
\begin{longtable}{@{}>{\RaggedRight\arraybackslash}p{4.0cm}>{\RaggedRight\arraybackslash}p{3.0cm}>{\RaggedRight\arraybackslash}p{2.4cm}>{\RaggedRight\arraybackslash}p{\dimexpr\textwidth-4.0cm-3.0cm-2.4cm-6\tabcolsep\relax}@{}}
\caption{Kontrak \emph{canonical artifact} \emph{knowledge graph}}\label{tab:rancangan-canonical-kg}\\
\toprule Kontrak/Elemen & \emph{Field} & Tipe & Fungsi\\ \midrule
\endfirsthead
\multicolumn{4}{l}{\footnotesize\emph{(lanjutan)}}\\ \toprule Kontrak/Elemen & \emph{Field} & Tipe & Fungsi\\ \midrule \endhead
\bottomrule \endfoot
\code{CanonicalGraph} & \code{document_id} & \emph{string} & Identitas dokumen asal \emph{graph}.\\
\code{CanonicalGraph} & \code{source_id} & \emph{string} & Identitas sumber \emph{graph}.\\
\code{CanonicalGraph} & \code{owned_node_ids} & \emph{array of} \code{string} & Daftar ID \emph{node} yang dimiliki \emph{profile}.\\
\code{CanonicalGraph} & \code{graph_builder} & \emph{PluginRef} & Identitas \emph{plugin} pembentuk \emph{graph}.\\
\code{CanonicalGraph} & \code{nodes} & \emph{array of} \code{Canonical}\allowbreak\code{Graph}\allowbreak\code{Node} & Daftar elemen \code{CanonicalGraphNode}.\\
\code{CanonicalGraph} & \code{edges} & \emph{array of} \code{Canonical}\allowbreak\code{Graph}\allowbreak\code{Edge} & Daftar elemen \code{CanonicalGraphEdge}.\\
\code{CanonicalGraph} & \code{metadata} & \emph{object} & \emph{Metadata} \emph{graph}.\\
\code{CanonicalGraphNode} & \code{id} & \emph{string} & Identitas simpul yang stabil.\\
\code{CanonicalGraphNode} & \code{type} & \emph{string} & Jenis simpul.\\
\code{CanonicalGraphNode} & \code{properties} & \emph{object} & Properti simpul.\\
\code{CanonicalGraphEdge} & \code{source} & \emph{string} & \emph{Endpoint} sumber relasi.\\
\code{CanonicalGraphEdge} & \code{target} & \emph{string} & \emph{Endpoint} tujuan relasi.\\
\code{CanonicalGraphEdge} & \code{type} & \emph{string} & Jenis relasi berarah.\\
\code{CanonicalGraphEdge} & \code{properties} & \emph{object} & Properti relasi.\\
\end{longtable}

\begin{longtable}{@{}p{3.2cm}p{3.2cm}p{\dimexpr\textwidth-3.2cm-3.2cm-4\tabcolsep\relax}@{}}
\caption{Konfigurasi \emph{plugin} \texttt{dense-semantic}}\label{tab:rancangan-config-indexer-dense}\\
\toprule \emph{Field} & Nilai \emph{default}/batas & Makna dan keluaran\\ \midrule
\endfirsthead
\multicolumn{3}{l}{\footnotesize\emph{(lanjutan)}}\\ \toprule \emph{Field} & Nilai \emph{default}/batas & Makna dan keluaran\\ \midrule \endhead
\bottomrule \endfoot
\code{indexed_metadata} & array, \emph{default} kosong, maksimum 32 item & \emph{Metadata} yang diproyeksikan ke dokumen indeks. Setiap item memiliki nama, path \emph{metadata}, dan tipe.\\
\code{embedding.provider} & enum local/openrouter/openai, \emph{default} \code{local} & \emph{Adapter} penyedia \emph{embedding}.\\
\code{embedding.model} & \emph{default} \code{BAAI/bge-m3} & Identitas model \emph{embedding}.\\
\code{embedding.api_key_ref} & opsional untuk \emph{provider} eksternal & Referensi kunci akses dari environment.\\
\code{embedding.api_key} & opsional untuk \emph{provider} eksternal & Nilai kunci akses inline.\\
\code{embedding.base_url} & \emph{default} URL \emph{provider} & Alamat \emph{service} \emph{embedding} eksternal.\\
\code{embedding.timeout_seconds} & number, \emph{default} 30 & Batas waktu permintaan \emph{embedding}.\\
\code{embedding.revision} & opsional & Revisi model \emph{embedding}.\\
\code{embedding.local_backend} & opsional & \emph{Backend} inferensi lokal.\\
\code{embedding.batch_size} & \emph{default} 32 & Jumlah teks dalam satu batch inferensi.\\
\code{embedding.device} & \emph{default} \code{auto} & Perangkat inferensi yang digunakan.\\
\end{longtable}

\endgroup
